\documentclass[a4paper]{article}
\usepackage[pdftex]{graphicx}
\usepackage[utf8]{inputenc}
\usepackage{enumerate}
\usepackage{icomma}
\usepackage{href-ul}
\hypersetup{
	colorlinks=true,
	linkcolor=blue,
	urlcolor=blue}
\usepackage{siunitx}
\sisetup{locale=DE} 
\usepackage{amssymb}
\usepackage{geometry}
\geometry{a4paper, top=15mm, left=15mm, right=15mm, bottom=15mm,
	headsep=10mm, footskip=12mm}

\begin{document}

\begin{enumerate}[1.]

%Aufgabe 1
{\bf \item Als Galileo Galilei die Geschwindigkeit des Lichts ermitteln wollte, war das nicht von Erfolg geprägt.}
\begin{enumerate}[a)]
\item Beschriebe stichpunktartig ein Messverfahren, mit welchem man die Lichtgeschwindigkeit bestimmen kann.
\item Gib das Symbol und den Wert der Lichtgeschwindigkeit an.
\end{enumerate}

%Aufgabe 2
{\bf \item Das Licht zweier Taschenlampen $L_1(1|2)$ und $L_2(2|6)$ trifft auf einen Körper,} der am Punkt $A(6,5|2)$ beginnt und im Punkt $B(6,5|4)$ endet. Hinter dem Körper befindet sich eine Wand, die die Strecke mit den Punkten $S_1(8|0)$ und $S_2(8|6)$ bildet. 

\begin{enumerate}[a)]
\item Zeichne die Punkte in ein Koordinatensystem und konstruiere den Schatten des Körpers
\item Teile den Schatten in verschiedene Bereiche ein und verwende dafür Fachbegriffe!
\end{enumerate}

%Aufgabe 3
{\bf \item Gib den Namen der gezeigten Linse an.}

\includegraphics[width=7 cm]{linsen100.pdf}

%Aufgabe 4
{\bf \item Das Keplersche und das Holländische (Galilei-)Fernrohr gehören wohl zu den bekanntesten. }
Nenne drei wesentliche Unterschiede zwischen diesen beiden.

%Aufgabe 5
{\bf \item Zur Verfügung steht eine Sammellinse mit der Brennweite $f=\SI{40}{\milli\meter}$.}

\begin{enumerate}[a)]
\item Konstruiere mit der Gegenstandsweite $g=\SI{15}{\milli\meter}$ und der Gegenstandsgröße 
$G=\SI{20}{\milli\meter}$ die Bildgröße $B$.
\item Gib an, welches Bild daraus entsteht und  nenne zwei Eigenschaften dieses Bildes.
\end{enumerate}

%Aufgabe 6
{\bf \item Erkläre, wie Kurzsichtigkeit zustande kommt und wie man sie beheben kann.}

%Aufgabe 7
{\bf \item Zwei der unsichtbaren Sonnenlichtbestandteile sind das ultraviolette Licht und das infrarote Licht.} Nenne drei Beispiele, wie man sich ultraviolette Strahlung zu Nutze macht.

%Aufgabe 8
{\bf \item Kreuze die zutreffenden Aussagen zur Dispersion an:}

$\square$ Die Auffächerung der Spektralfarben wird Dispersion genannt.\\
$\square$ Bei monochromatischem Licht entsteht stets ein Linienspektrum.\\
$\square$ UV--Licht und infrarotes Licht gehören nicht zu den Spektralfarben.\\ 
$\square$ Das weiße Sonnenlicht ist ein Beispiel für monochromatisches Licht.\\
$\square$ Eine durchgängige Auffächerung des Lichts nennt man kontinuierliches Spektrum.

%Aufgabe 9
{\bf \item Aufgabe}\\
\begin{minipage}{0,5\textwidth}
\begin{enumerate}[a)]
\item Erkläre stichpunktartig, was man unter Totalreflexion versteht.
\item Zeichne den Strahlengang durch das Prisma ein.
\end{enumerate}
\end{minipage}
\hfill
\begin{minipage}{0,45\textwidth}
\includegraphics[width=5 cm]{prisma100.pdf}
\end{minipage}

%Aufgabe 10
{\bf \item Eine fest eingespannte Sammellinse erzeugt von einem leuchtenden Gegenstand auf einem Milchglasschirm ein reelles Bild.} Was muss man alles verändern, wenn man ein kleineres, wiederum scharfes Bild des Gegenstandes erhalten will ?
\end{enumerate} 

\fbox{
	\begin{minipage}{0.5\textwidth}
		Weitere Arbeitsblätter gibt es unter 
		
		\href{https://www.okuyakl.de}{www.okuyakl.de}
	\end{minipage}
	\hfill
	\begin{minipage}{0.4\textwidth}
		\includegraphics[width=1.5 cm]{../../viecher/zwe03}
		\includegraphics[width=2 cm]{../../viecher/afanticon1}
		
\end{minipage}}

\end{document}

